\SourceLeaf{301}%
% A544-P0301-D01
\[
\begin{aligned}
(1)=(2)=(6)=(7)=(11)=(12)&=n\\
(3)=(10)&=\tfrac12(2n+1+t)\\
(4)=(9)&=\tfrac12(2n-t+u)\\
(5)=(8)&=\tfrac12(2n+1-u)
\end{aligned}
\]
\[
\begin{array}{r|r|r|r|r|r|r|r}
p&n&t&u&(1)&(3)&(4)&(5)\\\hline
5&0&1&1&0&1&0&0\\
29&1&3&3&1&3&1&0\\
53&2&3&5&2&4&3&0\\
101&4&7&5&4&8&3&2\\
149&6&7&7&6&10&6&3\\
173&7&7&9&7&11&8&3\\
197&8&5&11&8&11&11&3\\
269&11&11&7&11&17&9&8
\end{array}
\]

% A544-P0301-D02
Vierter Fall; $p=24n+17$.
\[
\begin{aligned}
2t&\quad\text{Anzahl der Classen für den Determinans }-p;\\
2u&\quad\text{Anzahl der Classen für den Determinans }-3p.
\end{aligned}
\]
\[
\begin{aligned}
(1)=(2)=(6)=(7)=(11)=(12)&=\tfrac16(6n+3+u)\\
(3)=(10)&=\tfrac16(6n+6+3t-2u)\\
(4)=(9)&=\tfrac16(6n+3-3t+u)\\
(5)=(8)&=\tfrac16(6n+6-2u)
\end{aligned}
\]
\[
\begin{array}{r|r|r|r|r|r|r|r}
p&n&t&u&(1)&(3)&(4)&(5)\\\hline
17&0&2&3&1&1&0&0\\
41&1&4&3&2&3&0&1\\
89&3&6&3&4&6&1&3\\
113&4&4&9&6&4&4&2\\
137&5&4&9&7&5&5&3\\
233&9&6&15&12&8&9&5\\
257&10&8&9&12&12&8&8
\end{array}
\]
\BookRule{.30\textwidth}

\SourceLeaf{302}%
% A544-P0302-D01
\par\bigskip\pdfbookmark[0]{Bemerkungen zu De nexu}{nexus-bemerkungen}
\begin{center}
\textsc{Bemerkungen zur Abhandlung}\\[1em]
\textsc{De nexu inter multitudinem classium, in quas formae binariae secundi gradus distribuuntur, earumque determinantem.}
\end{center}
\BookRule{.16\textwidth}
\par\medskip\pdfbookmark[1]{Zu I. und II.}{nexus-note-I-II}\centerline{Zu I. und II.}\smallskip

% A544-P0302-D02
Die zweite Formel für die Anzahl der innerhalb des Kreises liegenden Puncte (I. art. 3 und II. art. 5) ergiebt sich aus der Betrachtung des in denselben eingeschriebenen Quadrates, dessen Seiten den Coordinatenaxen parallel sind; die Vergleichung beider Formeln führt zu dem auch arithmetisch leicht zu beweisenden Satze
\[
r'+r''+\ldots+r^{(q)}=qq+r^{(q+1)}+r^{(q+2)}+\ldots+r^{(r)}
\]
aus welchem sich wieder die Richtigkeit der ersten von den beiden folgenden Regeln ergiebt, die sich auf einem besondern Blatt vorfanden:

% A544-P0302-D03
„Auflösungen der Gleichung $xx+yy\leqq A$; formula
\[
1+4\surd A+4\surd\tfrac12 A+8\Sigma(\surd(A-nn)-n)
\]
wo bei jeder Wurzel der Bruch weggelassen und von $n=1$ bis $n=\surd A$ „(soll heissen $\surd\tfrac12 A$)“ summirt wird.

% A544-P0302-D04
Andre Formel
\[
1+4\{A-\tfrac A3+\tfrac A5-\tfrac A7+\tfrac A9-\tfrac A{11}\ldots\}
\]
wo bei jedem Theil der Bruch weggelassen.“

% A544-P0302-D05
Diese letztere Formel folgt aus dem später (I. art. 6) zur Anwendung kommenden Satze über die Anzahl aller verschiedenen Darstellungen einer bestimmten Zahl durch die Form $xx+yy$ (vergl. Disqq. Arithm. art. 182, Note), welcher leicht in den folgenden umgeformt werden kann: die Anzahl der verschiedenen Darstellungen einer positiven ganzen Zahl $m$ durch die Form $xx+yy$ ist $=4(a-b)$, wo $a,b$ die Anzahlen der Divisoren von $m$ bedeuten, welche resp. von der Form $4n+1,\ 4n+3$ sind. Aus der Vergleichung
\SourceLeaf{303}%
% A544-P0303-D01
dieser arithmetischen Formel mit der (in I. art. 5 oder II. art. 4) durch geometrische Betrachtungen gewonnenen mittlern Darstellungsanzahl erhält man leicht und in aller Strenge das bekannte Resultat
\[
\frac\pi4=1-\tfrac13+\tfrac15-\tfrac17+\ldots
\]
welches in der Abhandlung (I. art. 7) durch eine ähnliche Vergleichung, aber mit Hülfe unendlicher Producte abgeleitet wird.
\BookRule{.25\textwidth}

% A544-P0303-D02
\par\medskip\pdfbookmark[1]{Zu III und IV.}{nexus-note-III-IV}\centerline{Zu III und IV.}\smallskip
Ist $C$ der Complex aller positiven, nicht eigentlich-äquivalenten formae proprie primitivae von negativem Determinant $-D$, und legt man den Variabeln dieser Formen je zwei Werthe bei, welche relative Primzahlen zu einander sind, so ist die Anzahl aller Darstellungen einer positiven ganzen Zahl $m$ gleich $\epsilon\psi(m)$, wo $\epsilon$ die Anzahl der Auflösungen der Gleichung $tt+Duu=1$, und $\psi(m)$ die Anzahl derjenigen Wurzeln $n$ der Congruenz $nn+D\equiv0\pmod m$ bedeutet, für welche die drei Zahlen $m,\ 2n$ und $\frac{nn+D}m$ ohne gemeinschaftlichen Divisor sind (Disqq. Arithm. art. 180). Der Factor $\epsilon$ ist $=4$ für $D=1$, in allen andern Fällen $=2$. Ist ferner $m=p^\pi p'^{\pi'}p''^{\pi''}\ldots$, wo $p,p',p''\ldots$ von einander verschiedene Primzahlen bedeuten, so ist $\psi(m)=\psi(p^\pi)\psi(p'^{\pi'})\psi(p''^{\pi''})\ldots$; bedeutet $\mathfrak A(m)$ die Anzahl der Wurzeln $n$ der Congruenz $nn+D\equiv0\pmod m$, und bedient man sich des von \textsc{Legendre} eingeführten, von \textsc{Jacobi} verallgemeinerten Zeichens, so ist $\psi(p^\pi)=\mathfrak A(p^\pi)=1+\left(\frac{-D}p\right)$, wenn $p$ nicht in $2D$ aufgeht, sonst aber $=\mathfrak A(p^\pi-\frac1p\mathfrak A(p^{\pi+1})$; die Anzahl $\mathfrak A(p^\pi)$ lässt sich immer leicht bestimmen (Disqq. Arithm. art. 104), für die Folge reicht aber die Bemerkung aus, dass $\mathfrak A(p^\pi)$ immer von $\pi$ unabhängig wird, sobald $\pi$ eine gewisse Grösse überschreitet.

% A544-P0303-D03
Legt man den Variabeln der in dem Complex $C$ enthaltenen Formen alle ganzzahligen Werthe ohne Ausnahme bei (Disqq. Arithm. art. 181), so wird die Anzahl $(m)$ aller Darstellungen der Zahl $m$ gleich $\epsilon f(m)$, wo $f(m)=\Sigma\psi\left(\frac m{\mu\mu}\right)$ ist, und das Summenzeichen sich auf alle quadratischen Divisoren $\mu\mu$ der Zahl $m$ bezieht. Hieraus folgt unmittelbar
\[
f(m)=f(p^\pi p'^{\pi'}p''^{\pi''}\ldots)=f(p^\pi)f(p'^{\pi'})f(p''^{\pi''})\ldots
\]
und
\[
f(p^\pi)=\psi(p^\pi)+\psi(p^{\pi-2})+\psi(p^{\pi-4})+\ldots
\]
welche Reihe so lange fortzusetzen ist, als die Exponenten $\pi,\ \pi-2,\ \pi-4\ldots$ nicht negativ werden. Wenn $p$ nicht in $2D$ aufgeht, so folgt hieraus
\[
f(p^\pi)=1+\left(\frac{-D}p\right)+\left(\frac{-D}{pp}\right)+\ldots+\left(\frac{-D}{p^\pi}\right)
\]
und allgemein, wenn $m$ relative Primzahl zu $2D$ ist,
\[
f(m)=\Sigma\left(\frac{-D}n\right)
\]
wo das Summenzeichen sich auf alle Divisoren $n$ der Zahl $m$ bezieht.

\SourceLeaf{304}%
% A544-P0304-D01
Aus diesen Bemerkungen ergiebt sich unmittelbar die Richtigkeit der im Text (III, 1, 2, 3) aufgestellten Sätze über die Anzahl $(m)$, wenn man für den ersten derselben noch die Bedingung hinzufügt, dass $D$ nicht durch $pp$ theilbar sein darf (die Bestimmung der Classenanzahl ist schon in den Disqq. Arithm. art. 256 auf den Fall zurückgeführt, in welchem $D$ durch kein Quadrat theilbar ist). Zugleich findet man, auch ohne Rücksicht auf diese Beschränkung, dass die unendliche Reihe
\[
1+\frac{f(p)}p+\frac{f(pp)}{pp}+\frac{f(p^3)}{p^3}+\ldots
\]
den Werth
\[
\frac1{1-\frac1p}\quad\text{oder}\quad\frac1{1-.\frac1p}\pdot\frac1{1-\left(\frac{-D}p\right)\frac1p}
\]
hat, je nachdem $2D$ durch die Primzahl $p$ theilbar oder nicht theilbar ist.
\BookRule{.25\textwidth}

% A544-P0304-D02
\par\medskip\pdfbookmark[1]{Zu V.}{nexus-note-V}\centerline{Zu V.}\smallskip
Die zu der Formel III hinzugefügte Bemerkung giebt den Weg an, auf welchem der Verf. zur Bestimmung der Anzahl $k$ der in dem Complex $C$ enthaltenen Formen gelangt ist. Aus geometrischen Betrachtungen (vergl. I. art. 5 und II. art. 4) ergiebt sich, dass der Grenzwerth, welchem sich der Quotient
\[
\frac{(1)+(2)+(3)+\ldots+(m)}m
\]
mit unbegrenzt wachsendem $m$ nähert, d. h. die mittlere Anzahl der Darstellungen einer unbestimmten positiven ganzen Zahl
\[
=k\frac\pi{\surd D}
\]
ist; ein zweiter Ausdruck für denselben Grenzwerth lässt sich auf verschiedene Arten aus der Natur der im Vorhergehenden bestimmten Anzahl $(m)=\epsilon f(m)$ der Darstellungen der Zahl $m$ ableiten. Der zu diesem Zweck von dem Verf. zunächst eingeschlagene Weg scheint nach den vorhandenen Bruchstücken (I. artt. 7, 8; III und IV) folgender gewesen zu sein.

% A544-P0304-D03
Ist $\theta(m)$ irgend eine Function der positiven ganzen Zahl $m$, und $p$ irgend eine Primzahl, so kann man aus $\theta(m)$ immer eine neue Function $\theta'(m)$ ableiten, deren Werth unabhängig davon ist, ob und wie oft $p$ als Factor in $m$ enthalten ist, und welche für alle durch $p$ nicht theilbaren Zahlen $m$ mit $\theta(m)$ übereinstimmt; eine solche Function erhält man, wenn man $\theta'(m)=\theta\left(\frac m{p^\pi}\right)$ setzt, wo $p^\pi$ die höchste in $m$ aufgehende Potenz von $p$ bedeutet; und man kann sagen, dass die Function $\theta'(m)$ aus $\theta(m)$ durch Elimination der Primzahl $p$ entsteht. Bildet man auf diese Weise aus $f(m)$ eine neue Function $f'(m)$ durch Elimination der Primzahl $2$, aus dieser die Function $f''(m)$ durch Elimination von $3$ u. s. f., so wird jede folgende dieser Functionen einen regelmässigern Verlauf haben, als die vorhergehenden; eliminirt man eine Primzahl nach der andern, wie sie ihrer Grösse nach auf einander folgen, so wird eine solche Function
\SourceLeaf{305}%
% A544-P0305-D01
$\theta(m)$ für unendlich viele Werthe von $m$ den Werth $f(1)=1$ haben, und namentlich für alle diejenigen Werthe von $m$, welche kleiner sind als die zuletzt eliminirte Primzahl. Durch unendliche Fortsetzung dieses Processes nähert man sich immer mehr der Function $f^\infty(m)$, welche für alle Werthe von $m$ den Werth $1$ hat, und deren mittlerer Werth folglich ebenfalls $=1$ ist. Gelingt es nun den mittlern Werth irgend einer Function $\theta(m)$ durch denjenigen der nächstfolgenden $\theta'(m)$ auszudrücken, so wird man auch den mittlern Werth der Function $f(m)$ durch eine unendliche Kette von Operationen finden können.

% A544-P0305-D02
Ist $p$ die Primzahl, durch deren Elimination $\theta'(m)$ aus $\theta(m)$ entsteht, so ist $\theta(m)=\theta'(m)f(p^\pi)$, wenn $p^\pi$ wieder die höchste in $m$ aufgehende Potenz von $p$ bedeutet. Für den Fall, dass $p$ nicht in $2D$ aufgeht, findet man hieraus leicht, dass
\[
\theta'(m)=\theta(mp)-\left(\frac{-D}p\right)\theta(m)
\]
ist; setzt man zur Abkürzung
\[
\begin{aligned}
\Theta(m)&=\theta(1)+\theta(2)+\ldots+\theta(m)\\
\Theta'(m)&=\theta'(1)+\theta'(2)+\ldots+\theta(m)
\end{aligned}
\]
so ergiebt sich
\[
\Theta'(mp)=\Theta(mp)-\left(\frac{-D}p\right)\Theta(m)
\]
und hieraus, wenn man mit $\omega,\omega'$ resp. die mittlern Werthe der Functionen $\theta(m),\theta'(m)$ bezeichnet,
\[
\omega=\frac{\omega'}{1-\left(\frac{-D}p\right)\frac1p}
\]

% A544-P0305-D03
Wenn aber die Primzahl $p$ in $2D$ aufgeht, so findet zwar zwischen den Functionen $\theta(m)$ und $\theta'(m)$ im Allgemeinen keine so einfache Beziehung mehr Statt; indessen ergiebt sich auf ähnliche Art leicht, dass in diesem Fall $\omega=\omega'$ ist. Ein anderer Weg, die Beziehung zwischen $\omega$ und $\omega'$ in beiden Fällen abzuleiten, ist folgender. Setzt man
\[
\vartheta(m)=\Sigma\theta(\mu)
\]
wo das Summenzeichen sich auf alle Zahlen $\mu$ bezieht, die nicht durch $p$ theilbar und ausserdem nicht grösser als $m$ sind, und bezeichnet man mit $m',m'',m'''\ldots$ resp. die grössten in $\frac mp,\frac{m'}p,\frac{m''}p\ldots$ enthaltenen ganzen Zahlen, so ist
\[
\begin{aligned}
\Theta(m)&=\vartheta(m)+\vartheta(m')f(p)+\vartheta(m'')f(pp)+\vartheta(m''')f(p^3)+\ldots\\
\Theta'(m)&=\vartheta(m)+\vartheta(m')+\vartheta(m'')+\vartheta(m''')+\ldots
\end{aligned}
\]
und hieraus folgt
\[
\frac\omega{\omega'}=(1-\tfrac1p)\{1+\frac{f(p)}p+\frac{f(pp)}{pp}+\frac{f(p^3)}{p^8}+\ldots\}
\]
was mit dem oben gefundenen Resultat übereinstimmt (vergl. die Note zu III und IV).

% A544-P0305-D04
Der mittlere Werth der Function $f(m)$ ist daher gleich dem unendlichen Product
\[
\Pi\frac1{1-\left(\frac{-D}p\right)\frac1p}
\]
